\section{Vertebrates to Mammals --- The Ancient Machinery of Choice}\label{05_vertebrates-vertebrates-to-mammals-the-ancient-machinery-of-choice}

The architecture by which a human brain selects what to do next is, at the level of cell types, neurotransmitters, receptors, ion channels, and circuit logic, \textbf{virtually identical to the architecture by which a lamprey selected what to do next more than 500 million years ago}. This is not a metaphor or a loose analogy. It is the empirically demonstrated conclusion of three decades of comparative neuroanatomy, and it changes what one can intelligibly mean by human agency. The basal ganglia and their dopaminergic modulators are not, as popular neuroscience persistently claims, a ``reptilian'' relic overruled by a more recent neocortex. They are the original action-selection engine of vertebrates, present and operational before jaws, before limbs, before anything resembling a cerebral cortex. Everything that came later --- the elaborated pallium, the six-layered neocortex of mammals, the granular prefrontal cortex of primates --- was built on top of and \emph{through} that engine, not as its replacement but as a way to feed it richer, more abstract, more temporally extended content. The chapter that follows traces this evolutionary lineage from cyclostomes through early mammals, with one central claim: human agency is the latest expression of an extraordinarily conserved biological computation, and its material substrate is therefore directly inspectable, directly damageable, and directly continuous with that of every vertebrate alive today.

\subsection{The vertebrate bauplan was set before the Cambrian closed}\label{05_vertebrates-the-vertebrate-bauplan-was-set-before-the-cambrian-closed}

By the late Cambrian, more than 500 million years ago, the vertebrate central nervous system had already crystallized into the three vesicles that remain the universal scaffold of every vertebrate brain examined to date: \textbf{prosencephalon (forebrain), mesencephalon (midbrain), and rhombencephalon (hindbrain)}. Lampreys, hagfishes, sharks, teleosts, frogs, lizards, birds, and mammals are all variations on this ground plan. The hindbrain houses cranial nerve nuclei, respiratory and cardiovascular pattern generators, vestibular nuclei, and the central pattern generators for locomotion. Its segmental organization into rhombomeres, each carrying a distinct combinatorial Hox code, is conserved from lamprey to mouse \href{https://openaccesspub.org/evolutionary-science/article/1641}{Open Access Pub} --- disrupting Hox, Pbx, or Meis function produces homeotic transformations of rhombomeres \href{https://www.ncbi.nlm.nih.gov/pmc/articles/PMC4489388/}{nih} in fish and mammals alike. The cerebellum, present in all jawed vertebrates \href{https://pmc.ncbi.nlm.nih.gov/articles/PMC3389513/}{PubMed Central} and rudimentary in cyclostomes, runs the same canonical circuit (mossy and parallel fibers, Purkinje cells, deep nuclei, climbing-fiber error signals from the inferior olive) wherever it appears, implementing supervised motor learning and forward-model prediction even in larval zebrafish. \href{https://www.science.org/doi/10.1126/sciadv.adi6470}{Science}

The midbrain divides into a dorsal \textbf{tectum} and a ventral \textbf{tegmentum}. The optic tectum, which mammalian anatomists call the superior colliculus, is one of the deepest sensorimotor integrators in the vertebrate brain: laminated, retinotopic, and aligned with a deeper motor map for orienting \href{https://link.springer.com/article/10.1007/s12264-022-00858-1}{Springer} and evasive movements in every species examined. The same brainstem-projecting tectal cell classes mediate the looming-evoked switch between approach and escape in lamprey, zebrafish, and mouse. \href{https://pmc.ncbi.nlm.nih.gov/articles/PMC6660747/}{PubMed Central} The tegmentum houses the dopaminergic \textbf{substantia nigra pars compacta (SNc)} and \textbf{ventral tegmental area (VTA)}, the GABAergic substantia nigra pars reticulata (SNr), and the mesencephalic locomotor region --- all conserved from cyclostomes to mammals.

The forebrain divides into diencephalon (thalamus, hypothalamus, epithalamus, pretectum) and telencephalon (pallium, subpallium). Lampreys possess each of these subdivisions, including a dorsal pallium with discrete retinotopic visual, somatosensory, and motor areas, \href{https://www.semanticscholar.org/paper/The-Basal-Ganglia-Over-500-Million-Years-Grillner-Robertson/97ee475166bf993d3e74ef14e056b2aab025e04d}{Semantic Scholar} \href{https://www.karger.com/Article/FullText/517492}{Karger} and a fully differentiated subpallial basal ganglia. The cross-vertebrate persistence of these subdivisions is rationalized by the \textbf{prosomeric model} of Puelles and Rubenstein, in which the forebrain is partitioned into transverse prosomeres and longitudinal zones defined by conserved patterning genes (Shh, Wnt, Otx, Gbx2, Pax6, Nkx2.1, Dlx, Emx, Lhx, Foxg1). Equivalent progenitor domains can be molecularly identified in lamprey, shark, teleost, frog, chick, and mouse brains. The bauplan is not a metaphor; it is a developmental fact written into shared genomes.

\subsection{Why ``the reptilian brain'' is one of neuroscience's most successful errors}\label{05_vertebrates-why-the-reptilian-brain-is-one-of-neurosciences-most-successful-errors}

Paul MacLean's \textbf{triune brain} model, formalized in his 1990 book \emph{The Triune Brain in Evolution}, parsed the human brain into three evolutionary tiers: a ``reptilian complex'' of basal ganglia and brainstem governing instinct, a ``paleomammalian'' limbic system governing emotion, and a ``neomammalian'' neocortex governing reason. The model became wildly influential --- adopted by Carl Sagan \href{https://www.msn.com/en-us/news/technology/seven-and-a-half-lessons-about-the-brain-by-lisa-feldman-barrett-review/ar-BB1eca70}{MSN} in \emph{The Dragons of Eden}, repeated in countless textbooks, and still invoked today in management seminars and meditation apps. It is also empirically false. As the comparative-neurobiology review \textbf{``Your Brain Is Not an Onion With a Tiny Reptile Inside'' (Cesario, Johnson \& Eisthen, 2020)} put the consensus, the triune scheme is ``one of the most successful and widespread errors in all of science.'' \href{https://cpbotha.github.io/2021/06/05/book-notes-seven-and-a-half-lessons-about-the-brain/}{Cpbotha}

Four objections matter for the argument of this chapter. First, \textbf{the so-called reptilian basal ganglia long predate reptiles}. Lampreys, which diverged from the gnathostome lineage more than 500 million years ago, \href{https://www.sciencedirect.com/science/article/pii/S0960982216306807}{ScienceDirect} possess a fully differentiated basal ganglia \href{https://pubmed.ncbi.nlm.nih.gov/21700460/}{PubMed} with the same direct/indirect pathways, \href{https://pubmed.ncbi.nlm.nih.gov/23318875/}{PubMed} the same D1/D2 dopamine signaling, and the same tonic GABAergic output as a rodent. \href{https://pmc.ncbi.nlm.nih.gov/articles/PMC3853485/}{PubMed Central} \href{https://www.sciencedirect.com/science/article/pii/S0960982214012998}{ScienceDirect} Calling this circuitry ``reptilian'' is a chronological misnomer of more than 250 million years. Second, reptiles, birds, and fishes themselves possess pallial homologues that support sophisticated cognition; the avian pallium runs tool use, episodic-like memory, and metacognition without any neocortex at all. \href{https://ucrisportal.univie.ac.at/en/publications/cognition-without-cortex/}{University of Vienna} Third, the supposed ``paleomammalian'' limbic structures --- hippocampus, amygdala, hypothalamus, septum, habenula --- are present in lampreys, fishes, amphibians, reptiles, and birds. Fourth, and most fundamentally, \textbf{brains do not stack; they elaborate}. Comparative connectomics, shared developmental codes, and molecular phylogenetics show that vertebrate brain evolution proceeds by changes in proportional size, neuronal number, and elaboration of pre-existing fields, \href{https://global.oup.com/ushe/product/principles-of-brain-evolution-9780878938209}{Oxford University Press} not by accretion of new layers atop older ones. The triune model is a modern incarnation of the discredited \emph{scala naturae}, the Aristotelian ladder of progress culminating in humans. Living lampreys, frogs, lizards, crows, and humans are each terminal twigs of independent lineages, none ancestral to another.

This matters because the triune story underwrites a specific, widely held picture of agency: that ``rational choice'' is implemented in the cortex while ``instinct'' is implemented in older subcortical structures, and that willpower is the cortex governing the brainstem. \textbf{The actual architecture is the opposite}. The cortex feeds proposals into a far older selection engine, which decides what actually happens.

\subsection{The basal ganglia: a 500-million-year-old solution to the selection problem}\label{05_vertebrates-the-basal-ganglia-a-500-million-year-old-solution-to-the-selection-problem}

In 1999 Peter Redgrave, Tony Prescott, and Kevin Gurney articulated the framing that has organized basal-ganglia neuroscience ever since. The vertebrate body has only one set of muscles. Every cortical, limbic, and brainstem command system competes for access to that final common motor pathway. \textbf{A selection problem arises whenever multiple parallel processes seek simultaneous access to a restricted resource}, and any organism that survives must implement an arbitration mechanism. Redgrave and colleagues proposed that the vertebrate basal ganglia evolved precisely as a centralized selection device --- a switching circuit specialized to resolve conflicts over access to limited motor and cognitive resources. Three features of basal-ganglia organization fit this hypothesis perfectly. The basal ganglia receive convergent input from virtually every cortical and limbic command source. Their output is uniformly inhibitory and tonically active, \href{https://pubmed.ncbi.nlm.nih.gov/23318875/}{PubMed} so that downstream targets are \emph{off by default}. And selection is implemented by \textbf{focal disinhibition}: a winning channel transiently silences its corresponding output neurons, releasing one motor program while losing channels remain (or are made even more) suppressed. \href{https://eprints.whiterose.ac.uk/id/eprint/107033/1/Redgrave\%20Neuroscience\%201999\%20preprint.pdf}{White Rose Research Online}

The mammalian basal ganglia comprise five core nuclei. The \textbf{striatum} is the principal input, divided into a dorsal portion (caudate and putamen) and a ventral portion (nucleus accumbens). About 95\% of striatal cells are GABAergic medium spiny projection neurons (SPNs), with sparse cholinergic and parvalbumin-positive interneurons. The \textbf{globus pallidus}, divided into external (GPe) and internal (GPi) segments, receives striatal input and contributes the principal output. The \textbf{subthalamic nucleus (STN)} is the only glutamatergic basal-ganglia nucleus; it receives input from GPe and from cortex directly (the ``hyperdirect'' pathway) and excites the output nuclei. \href{https://pmc.ncbi.nlm.nih.gov/articles/PMC4771745/}{PubMed Central} The \textbf{substantia nigra}, divided into pars compacta (SNc, dopaminergic) and pars reticulata (SNr, GABAergic output), supplies dopamine to the striatum and shares the output role with GPi. The general operating rule is decisive: \textbf{GPi and SNr together hold downstream motor centers under continuous, high-frequency tonic inhibition at rest}. \href{https://www.eurekaselect.com/article/133535}{Bentham Science} \href{https://www.researchgate.net/publication/228035209_Anatomy_of_the_Corpus_Striatum_and_Brain_Stem_Integrating_Systems}{ResearchGate} Nothing happens by default.

Action selection runs on two opposing pathways through this circuitry. The \textbf{direct (``Go'') pathway} consists of striatal SPNs expressing D1 dopamine receptors plus substance P and dynorphin. These cells project monosynaptically to GPi and SNr. \href{https://www.sciencedirect.com/science/article/pii/S0960982216306807}{ScienceDirect} \href{https://pmc.ncbi.nlm.nih.gov/articles/PMC8710883/}{PubMed Central} When cortical glutamatergic input excites a specific population of D1 SPNs, those SPNs powerfully inhibit their corresponding GPi/SNr neurons, producing a transient pause in tonic firing; the downstream motor target is disinhibited, and the program runs. \href{https://www.sciencedirect.com/science/article/pii/S0960982214012998}{ScienceDirect} The \textbf{indirect (``No-Go'') pathway} consists of striatal SPNs expressing D2 receptors and enkephalin. \href{https://pubmed.ncbi.nlm.nih.gov/23318875/}{PubMed} These cells project to GPe; \href{https://pmc.ncbi.nlm.nih.gov/articles/PMC8710883/}{PubMed Central} GPe normally inhibits STN; D2 SPN activation therefore disinhibits STN, which then drives GPi/SNr more strongly, \emph{deepening} tonic inhibition and suppressing competing actions. Dopamine, acting via D1 receptors (Gs-coupled, increasing cAMP), enhances D1-SPN excitability and corticostriatal long-term potentiation, favoring the selected action; dopamine acting via D2 receptors (Gi/o-coupled, decreasing cAMP) suppresses D2-SPN excitability, releasing the brake on indirect drive. The push-pull architecture is now elegantly confirmed by optogenetics: stimulating D1 SPNs in mice produces movement, while stimulating D2 SPNs produces a Parkinsonian freeze.

The decisive demonstration that \textbf{this entire architecture predates the divergence of jawless and jawed vertebrates} comes from the Karolinska laboratory of Sten Grillner, working with Brita Robertson and Marcus Stephenson-Jones. Their 2011 paper in \emph{Current Biology}, ``Evolutionary conservation of the basal ganglia as a common vertebrate mechanism for action selection,'' used immunohistochemistry, tract tracing, and whole-cell patch-clamp recordings in \emph{Lampetra fluviatilis} to show that the lamprey possesses a striatum with mammal-type SPNs (inwardly rectifying, plateau-generating), substance-P-positive direct-pathway neurons projecting to a GPi homolog, enkephalin-positive indirect-pathway neurons projecting to a GPe homolog, a glutamatergic STN receiving GPe input and projecting to GPi, and a tonically active GABAergic GPi inhibiting brainstem motor centers. \href{https://www.sciencedirect.com/science/article/pii/S0960982214012998}{ScienceDirect} A 2012 paper in the \emph{Journal of Comparative Neurology} established the lamprey SNr homolog and its dual-output role with GPi. \href{https://onlinelibrary.wiley.com/doi/abs/10.1002/cne.23087}{Wiley Online Library} A 2013 paper in \emph{PNAS} established the lamprey homologue of the habenular circuit that controls aversive value signaling to dopamine and serotonin systems. A 2013 paper in the \emph{Journal of Neuroscience} by Ericsson and colleagues showed in lamprey the \emph{opposing} D1/D2 modulation of direct- and indirect-pathway SPNs --- dopamine excites D1 cells and inhibits D2 cells in lamprey exactly as in mammal. As the keystone 2016 review by Grillner and Robertson, \textbf{``The Basal Ganglia Over 500 Million Years''}, concluded: ``Transmitters, connectivity and membrane properties are virtually identical in lamprey and rodent basal ganglia.'' \href{https://www.sciencedirect.com/science/article/pii/S0960982216306807}{ScienceDirect} \href{https://www.researchgate.net/publication/228035209_Anatomy_of_the_Corpus_Striatum_and_Brain_Stem_Integrating_Systems}{ResearchGate}

The same conservation extends laterally across the gnathostome radiation. Reiner, Medina, and Veenman documented direct/indirect pathway architecture in fish, amphibians, reptiles, and birds. \href{https://www.sciencedirect.com/science/article/abs/pii/S0165017398000162}{ScienceDirect} The 2004 reclassification of the avian telencephalon (discussed below) corrected a century-old error and recognized that bird ``paleostriatum'' and ``archistriatum'' structures are true homologs of mammalian striatum and pallidum, with songbird Area X running the same direct/indirect logic for vocal learning. \textbf{What evolution added in mammals was scale and modular replication, not redesign.} The number of selection modules multiplied --- motor, oculomotor, \href{https://www.sciencedirect.com/science/article/pii/S0960982216306807}{ScienceDirect} dorsolateral prefrontal, lateral orbitofrontal, and limbic loops, each performing the same canonical disinhibitory computation on different content. Cortical input to striatum massively expanded. But the engine itself is a 500-million-year-old inheritance.

\subsection{Dopamine: the ancient currency of value and vigor}\label{05_vertebrates-dopamine-the-ancient-currency-of-value-and-vigor}

The dopaminergic modulators of the basal ganglia are equally ancient. Lampreys possess a \textbf{posterior tuberculum} that corresponds anatomically and functionally to the mammalian SNc and VTA, \href{https://www.sciencedirect.com/science/chapter/handbook/abs/pii/B9780443298677000074}{ScienceDirect} with the same connectome --- afferents from striatum, lateral habenula, cortex/pallium, and pedunculopontine; efferents to striatum, optic tectum, and pallium. Pérez-Fernández and colleagues in 2014 traced this in detail; in 2017 they showed in \emph{Neuron} that single SNc neurons in lamprey send branches to \emph{both} striatum and optic tectum, \href{https://www.cell.com/neuron/fulltext/S0896-6273(17)30920-0}{Neuron} \href{https://mdpi.com/1422-0067/22/20/11284/htm}{MDPI} implementing salience-driven dopaminergic modulation across both selection levels simultaneously. The full Dahlström-Fuxe taxonomy of dopaminergic cell groups (A8 through A16, including the hypothalamic A11--A15 populations) applies across all vertebrates investigated, \href{https://en.wikipedia.org/wiki/Dopaminergic_cell_groups}{Wikipedia} including lamprey. Even pre-vertebrate chordates such as amphioxus and tunicates possess dopamine-positive forebrain neurons. \href{https://www.sciencedirect.com/science/chapter/handbook/abs/pii/B9780443298677000074}{ScienceDirect} \textbf{The dopaminergic system is older than vertebrates themselves.}

Receptor signaling is equally conserved. The lamprey D2 receptor cDNA, cloned by Robertson and colleagues in 2012, shows close phylogenetic relationship with vertebrate D2 receptors \href{https://www.sciencedirect.com/science/chapter/handbook/abs/pii/B9780443298677000074}{ScienceDirect} and almost 100\% identity in the transmembrane domains containing the dopamine-binding amino acids. D1-like receptors (D1A, D1B/D5) are present in all vertebrates from lampreys to mammals; D2, D3, and D4 receptors likewise. \href{https://www.sciencedirect.com/science/chapter/handbook/abs/pii/B9780443298677000074}{ScienceDirect} The signaling polarity --- D1 = Gs/cAMP↑ → excitation, D2 = Gi/o/cAMP↓ → inhibition --- is uniform across phyla. \href{https://www.sciencedirect.com/science/chapter/handbook/abs/pii/B9780443298677000074}{ScienceDirect}

Functionally, dopamine does in lamprey what it does in human. \textbf{Thompson, Ménard, Pombal and Grillner showed \href{https://www.semanticscholar.org/paper/The-Basal-Ganglia-Over-500-Million-Years-Grillner-Robertson/97ee475166bf993d3e74ef14e056b2aab025e04d}{Semantic Scholar} in 2008 that MPTP-induced forebrain dopamine depletion in lamprey reduces striatal dopamine to roughly 15\% of control values and produces dramatic reductions in spontaneous swimming and bout duration --- a Parkinson-like phenotype in a 560-million-year-old lineage.} The phasic burst signature of dopamine reward-prediction error described by Wolfram Schultz in primates is mirrored in the lamprey GPh→lateral habenula→SNc circuit characterized by Stephenson-Jones in 2012 and 2013, and extended to mice in his 2016 \emph{Nature} paper showing the same circuit encodes positive and negative reward prediction errors. \href{https://www.nature.com/articles/nature19845}{Nature} Dopamine-dependent corticostriatal LTP and LTD operate by the same D1-LTP/D2-LTD logic across vertebrates. Ventral striatal/accumbens dopamine drives appetitive approach in mammals; analogous circuitry handles valence-based decisions in lamprey.

The implications for the chapter's argument are direct. \textbf{Whatever incentive-driven action selection is, it is implemented by an ancient, highly conserved, dopamine-modulated gating circuit that vertebrates have used continuously for half a billion years.} When that circuit is damaged --- by Parkinson's disease, by chronic cocaine, by neuroleptic blockade, by decades of compulsive habit formation --- the pathology is not the corruption of some uniquely human faculty. It is the corruption of vertebrate action machinery, exposed by the breakdown.

\subsection{The thalamus closes the loop}\label{05_vertebrates-the-thalamus-closes-the-loop}

The thalamus is the great relay and integrative hub of the diencephalon, partitioned across all vertebrates into three molecularly defined territories: a \textbf{dorsal thalamus} (glutamatergic relay nuclei projecting to pallium --- LGN for vision, MGN for audition, VPL/VPM for somatosensation, anterior nuclei for limbic content, mediodorsal for prefrontal, pulvinar for higher-order vision), a \textbf{ventral thalamus} (the predominantly GABAergic prethalamus including zona incerta and the reticular thalamic nucleus), and an \textbf{epithalamus} (habenular complex and pineal organ). The thalamus already had this fundamental partitioning at the divergence of cyclostomes more than 500 million years ago, with the lamprey dorsal thalamus receiving direct retinal, tectal, and lemniscal input, projecting to medial pallium, and reciprocally connected with striatum and tectum.

What matters for the action-selection argument is that the thalamus is the obligatory return arm of the basal-ganglia loop. \textbf{GPi and SNr project GABAergically to motor and prefrontal thalamic nuclei (VA, VL, MD), which excite cortex; cortex re-targets striatum, closing the loop}, the architecture that Garrett Alexander, Mahlon DeLong, and Peter Strick formalized in their 1986 review of parallel cortico-striato-pallido-thalamo-cortical circuits. Crucially, this loop is conserved in lamprey: the lamprey GPi homolog and SNr project to dorsal thalamus and tectum, \href{https://www.sciencedirect.com/science/article/pii/S0960982214012998}{ScienceDirect} and lamprey thalamus reciprocally targets striatum and pallium. Within this conserved framework, mammals show striking expansions of specific dorsal thalamic territories --- the pulvinar enormously enlarged in primates, the mediodorsal nucleus co-expanded with prefrontal cortex --- but these are quantitative elaborations on a conserved template, not novel additions.

\subsection{The pallium and the warning birds give us}\label{05_vertebrates-the-pallium-and-the-warning-birds-give-us}

The pallium is the dorsal telencephalon, and across all jawed vertebrates it subdivides into four molecularly defined progenitor domains: medial (MP), dorsal (DP), lateral (LP), and ventral pallium (VP), demarcated by Emx1, Emx2, Tbr1, Pax6, Lhx2, and Lhx9. In mammals these domains generate hippocampus (MP), neocortex (DP), olfactory and piriform cortex (LP), and claustro-amygdalar nuclei (VP). \href{https://pmc.ncbi.nlm.nih.gov/articles/PMC4429232/}{PubMed Central +2} In sauropsids the same fields generate medial cortex, dorsal cortex/hyperpallium, and the dorsal ventricular ridge. The pallial bauplan is shared; what differs is how its fields are folded, layered, or aggregated into nuclei.

The defining mammalian innovation is the \textbf{six-layered isocortex} built from dorsal pallium, with layer I molecular, layers II/III supragranular associative, layer IV granular thalamorecipient, layer V infragranular pyramidal output, and layer VI corticothalamic. Embedded in this sheet is the canonical microcircuit characterized by Douglas and Martin: thalamic afferents excite layer 4 → layer 4 to layers 2/3 → layers 2/3 to layers 5/6 output → layer 5 descends to brainstem and spinal cord while layer 6 closes the corticothalamic loop. Cranial endocast evidence (Rowe, Macrini, and Luo, \emph{Science} 2011) places major neocortical expansion in the early mammaliaforms \emph{Morganucodon} and \emph{Hadrocodium} around 200 million years ago. Two complementary hypotheses explain its origin: Aboitiz proposes that the six-layered cortex arose through expansion and reorganization of reptilian dorsal pallium, driven by upregulation of dorsal signaling centers (cortical hem, BMPs/Wnts), amplification of reelin-producing Cajal--Retzius cells, and an inside-out radial migration program. \href{https://pmc.ncbi.nlm.nih.gov/articles/PMC4604247/}{PubMed Central} \href{https://www.sciencedirect.com/science/article/pii/S0301008220301209}{ScienceDirect} Karten's complementary cell-type homology hypothesis, vindicated molecularly by Dugas-Ford, Rowell, and Ragsdale in 2012, argues that the principal cell types of mammalian neocortex (layer 4 thalamorecipient, layer 5 brainstem-projecting) have homologs as separate nuclei distributed across the sauropsid DVR and Wulst --- confirmed by selective expression of L4 markers (EAG2, RORB) and L5 markers (ER81/ETV1, FEZF2, PCP4) in the predicted DVR nuclei across chicken and zebra finch. \href{https://pmc.ncbi.nlm.nih.gov/articles/PMC3479531/}{PubMed Central} \href{https://www.researchgate.net/publication/231742550_Cell-type_homologies_and_the_origins_of_the_neocortex}{ResearchGate} The deep lesson, articulated by Briscoe and Ragsdale in \emph{Science} in 2018, is that \textbf{natural selection preserves the integrity of cell types and information-processing circuits while three-dimensional architectures are far less constrained}. \href{https://pmc.ncbi.nlm.nih.gov/articles/PMC11098553/}{PubMed Central}

The most striking confirmation comes from birds. For a century, avian forebrain regions carried Edinger-era names (paleostriatum, archistriatum, neostriatum, hyperstriatum) reflecting the mistaken belief that nearly the entire bird telencephalon was striatal. This nomenclature rationalized the assumption that birds were cognitively limited --- ``bird-brained.'' In 2004 the Avian Brain Nomenclature Consortium, led by Anton Reiner and Erich Jarvis, overturned the terminology entirely. They renamed the truly striatal regions and reclassified the bulk of the cerebrum as \textbf{pallium}: hyperpallium, mesopallium, nidopallium, and arcopallium. \href{https://dukespace.lib.duke.edu/items/1132155b-17c6-4fa1-a7ae-8909390582f8}{Duke University} \href{https://pmc.ncbi.nlm.nih.gov/articles/PMC2507884/}{PubMed Central} Sixteen years later, \textbf{Stacho and colleagues, working with Onur Güntürkün, demonstrated in \emph{Science} (2020) that the avian pallium contains iteratively repeated, column-like neuronal circuitry} orthogonal to nuclear boundaries, \href{https://www.science.org/doi/10.1126/science.abc5534}{Science} \href{https://science.sciencemag.org/content/369/6511/eabc5534}{Science} with thalamic input → granular-like sensory recipient → secondary processing → output projection --- the canonical mammalian cortical microcircuit, implemented without cortical layers. In the same issue Nieder and colleagues recorded single neurons in the carrion crow's nidopallium caudolaterale during a near-threshold detection task and found activity tracking the bird's \emph{perceptual report} --- an empirical neural correlate of sensory consciousness in a brain with no neocortex.

The behavioral evidence is just as decisive. Corvids manufacture and use tools, plan caches for the next day, recognize themselves in mirrors, attribute knowledge to competitors. Parrots demonstrate inference by exclusion. Olkowicz and colleagues showed in 2016 that parrots and corvids pack roughly twice the neurons into a brain of equivalent mass as primates do; a blue-and-yellow macaw possesses about 1.9 billion pallial neurons, exceeding many monkeys with much larger brains. The avian nidopallium caudolaterale serves as a functional analog of mammalian prefrontal cortex, mediating working memory, rule learning, response inhibition, and reward valuation --- but it derives from ventral pallium, not dorsal pallium, and is therefore not homologous to PFC at the level of embryonic territory.

The implication for any account of human cognition is sharp. \textbf{Complex cognition does not require a six-layered mammalian neocortex.} What it requires are dense populations of glutamatergic projection neurons organized into iteratively repeated input--processing--output motifs, embedded in thalamopallial loops, richly interconnected through associative hubs, and modulated by a conserved subcortical action-selection engine. Whether those motifs are stacked into six layers (mammals), partially segregated within three layers (turtles), or aggregated into discrete nuclei (birds) appears to matter less than that the motifs are present, scaled, and integrated. Any theory that locates human consciousness or agency \emph{specifically} in the unique architecture of mammalian neocortex is empirically refuted by ravens solving multistep tool problems and by carrion crows generating perceptual-report-related neuronal signaling. \textbf{Cognition is substrate-flexible but computation-constrained.} The pallium, not the neocortex, is the deep evolutionary substrate of complex cognition --- and the engine that converts cognition into action is the same vertebrate basal-ganglia circuit running underneath both lineages.

\subsection{Limbic structures: motivation is older than mammals}\label{05_vertebrates-limbic-structures-motivation-is-older-than-mammals}

The amygdala, hippocampus, and hypothalamus are the same kind of story. They are not mammalian inventions; they are elaborations of an ancient vertebrate plan, with homologues in lampreys, fishes, amphibians, reptiles, and birds.

The \textbf{amygdala} is histogenetically composite, arising from both pallial and subpallial origins to generate four canonical components: the basolateral complex (deep pallial), the cortical nucleus (superficial pallial, olfactory-receiving), the medial nucleus (subpallial, vomeronasal), and the central nucleus (subpallial, striatal in origin, autonomic output). Moreno and González established that this dual pallial/subpallial bauplan is shared across all tetrapods, including amphibians. In reptiles the basolateral pallial amygdala corresponds to the posterior dorsal ventricular ridge; in birds the posteromedial arcopallium is its homolog, and lesions of avian arcopallium abolish conditioned fear and active-avoidance learning. In teleost fish the medial zone of the dorsal telencephalon (Dm) is the pallial amygdala homolog, and Dm lesions selectively impair avoidance learning --- the same phenotype as basolateral amygdala lesions in mammals. \textbf{Joseph LeDoux's dual-pathway fear circuit} --- a fast subcortical ``low road'' from sensory thalamus directly to lateral amygdala, and a slower cortical ``high road'' --- has homologous configurations across rodents, primates, amphibians, and fish. Janak and Tye's 2015 \emph{Nature} review made the conservation explicit: ``The amygdala is a brain region that is important for emotional processing, the circuitry and function of which appears to be largely conserved across vertebrates.'' Modern circuit dissection has revealed the basolateral amygdala as a \textbf{valence processor}, not a fear center: distinct projection-defined populations encode positive and negative outcomes (BLA→nucleus accumbens for reward seeking; BLA→central amygdala for fear; BLA→ventral hippocampus for anxiety), and these projections feed directly into the basal-ganglia action-selection apparatus.

The \textbf{hippocampus} arises from medial pallium across all vertebrates. In mammals it implements both the cognitive map of O'Keefe and the Mosers (place cells, grid cells, head-direction cells, border cells, recognized by the 2014 Nobel Prize) and the relational/episodic memory of Eichenbaum. The medial cortex of reptiles is its homolog; lesions impair allocentric spatial learning. The avian hippocampal formation, lacking an obvious dentate gyrus, nonetheless supports place-related firing, food-caching memory in chickadees and jays, and large-scale homing navigation in pigeons; Bingman and colleagues have established hemispheric lateralization paralleling human hippocampal asymmetries. In teleost fish the lateral zone of the dorsal pallium (Dl) is the hippocampal homolog; goldfish Dl lesions selectively impair allocentric place learning while sparing cued strategies, an exact phenocopy of mammalian hippocampal lesions. Recent neurologger recordings have identified place-like, head-direction, border, and even grid-like cells in the goldfish and zebrafish lateral pallium. \textbf{The neural representation of space is at least 400 million years old}, and the elaboration of hippocampal architecture through reptiles to mammals --- the appearance of laminated CA1--CA3, the dentate gyrus, the dramatic expansion of cortical input --- reflects increasing capacity for temporal sequencing, autobiographical context, and episodic-like recall.

The \textbf{hypothalamus} is among the most conserved structures in the vertebrate brain. Genoarchitectonic studies using Nkx2.1, Otp, Shh, Dlx, and Sim1 demonstrate conserved preoptic, anterior, tuberal, posterior, and mammillary compartments from lampreys through mammals. Lampreys possess identifiable hypothalamic territories with comparable molecular identities and functional hypothalamic-pituitary-gonadal and -thyroid axes. The hypothalamic dopaminergic cell groups A11 through A15 --- diencephalospinal, tuberoinfundibular, zona incerta, periventricular, and anteroventral periventricular --- are present throughout gnathostomes. The lateral hypothalamic area is the canonical motivational hub: electrical or optogenetic stimulation elicits voracious feeding, drinking, and locomotor activation; lesions cause aphagia and adipsia. The lateral hypothalamus contains the \textbf{orexin/hypocretin} neurons whose projections to VTA and nucleus accumbens provide the principal hypothalamic drive on the mesolimbic dopamine system. Orexin homologs are present in birds, fish, and amphibians; orexin loss underlies the evolution of sleep loss in cavefish.

These structures together generate motivated, value-based action through an integrated loop with the basal ganglia and VTA. \textbf{Sensory input converges on the basolateral amygdala, which assigns Pavlovian valence; the hippocampus provides spatial, contextual, and temporal information; the hypothalamus signals current homeostatic need; and the VTA integrates these inputs to compute reward prediction errors and motivational signals via dopamine release in the nucleus accumbens and dorsal striatum}. The nucleus accumbens functions as the limbic-motor interface, routing motivational output to ventral pallidum and thence to the same disinhibition machinery that gates motor action. Kent Berridge's two decades of work has established that mesolimbic dopamine encodes \textbf{incentive salience (``wanting'')}, not hedonic pleasure (``liking''): dopamine-depleted rats remain aphagic yet show normal hedonic facial reactions to sucrose, while sensitization of mesolimbic dopamine in addiction amplifies wanting without proportionate increases in liking. Hedonic liking depends on small opioid and endocannabinoid hotspots in the nucleus accumbens shell and ventral pallidum; wanting depends on the broad mesolimbic dopamine projection. \textbf{The same conserved circuit that selects swimming bouts in a lamprey selects compulsive behaviors in a person with addiction.} It is the breakdown of the system, not its invention, that human pathology reveals.

\subsection{Synapsids to mammals: layering sophistication onto ancient machinery}\label{05_vertebrates-synapsids-to-mammals-layering-sophistication-onto-ancient-machinery}

The mammalian brain assembled gradually along the synapsid lineage that diverged from sauropsids \textasciitilde320 million years ago. Pelycosaurs gave way to non-cynodont therapsids, then cynodonts, then mammaliaforms such as \emph{Morganucodon} and \emph{Hadrocodium} in the Early Jurassic \textasciitilde200 million years ago, and finally crown Mammalia. Across that 120-million-year corridor, characters distinctive of mammals --- the secondary palate, heterodonty, single dentary, three middle-ear ossicles, endothermy, fur, lactation, the six-layered neocortex --- assembled mosaically. Cranial endocast evidence from Rowe, Macrini, and Luo's 2011 \emph{Science} paper on \emph{Morganucodon} and \emph{Hadrocodium} identified three sequential pulses of brain expansion: an initial wave centered on enlarged olfactory bulbs and piriform cortex plus a primitive somatosensory neocortex and cerebellum (\textasciitilde200 mya), a second olfactory expansion in the \emph{Hadrocodium} grade, and a third pulse at the origin of crown Mammalia coupled with full ossification of ethmoid turbinals. Walls's ``\textbf{nocturnal bottleneck}'' rationalizes the selective context: small, burrowing, insectivorous mammaliaforms operating in low-light niches under the diurnal dominance of dinosaurs would have placed strong selective pressure on olfaction, somatosensation through vibrissae, audition through the new three-bone middle ear, and integrative neural circuitry capable of binding these channels into action.

Within the conserved mammalian bauplan, neocortical expansion across mammalian radiations follows two scaling phenomena. The first is \textbf{encephalization}, captured by Jerison's encephalization quotient, which trends upward in carnivores, cetaceans, and especially primates; humans reach an EQ of about 7. The second is \textbf{tangential areal expansion}, formalized by Rakic's radial unit hypothesis: a small lengthening of the symmetric-division phase in ventricular zone progenitors multiplicatively expands the founder pool and the number of cortical radial units, generating the more than thousandfold variation in cortical surface area across mammals with only modest changes in thickness. Subventricular-zone amplification, especially the outer SVZ in primates, drives upper-layer neuron generation, dense corticocortical connectivity, and gyrification. \textbf{What this expansion adds is real-estate for parcellation} --- more abstract hierarchical representations, more multimodal integration zones, longer temporal integration windows. Anterior cortical areas in primates integrate information over seconds and longer; posterior primary sensory areas integrate over milliseconds. The hierarchical extension of timescales is the central computational consequence of mammalian cortical scaling.

The \textbf{prefrontal cortex} is the apex of this hierarchy. Defined by dense projections from the mediodorsal thalamic nucleus, the PFC comprises agranular, dysgranular, and granular cytoarchitectonic regions. Standard subdivisions include dorsolateral PFC (areas 9, 46, 8) for working memory and rule maintenance, ventrolateral PFC for stimulus-response selection and language, orbitofrontal cortex for value representation, medial PFC for self-referential cognition, anterior cingulate cortex for conflict and effort integration, and frontopolar cortex (area 10) for prospective and meta-cognitive function. As Preuss and Wise synthesized in 2022, agranular medial frontal and orbital cortex are shared by all therian mammals, but \textbf{granular PFC, with a well-developed cell-rich layer 4, is a primate innovation}, evolving in early primates or in the last common ancestor of primates and tree shrews. Additional granular fields --- the dorsolateral PFC, dorsomedial PFC, frontal eye fields, granular OFC --- emerged in the simian (anthropoid) lineage, with humans further expanding granular PFC and its long-range corticocortical connectivity.

Functionally, Patricia Goldman-Rakic established that DLPFC pyramidal neurons exhibit \textbf{persistent delay-period activity} maintaining task-relevant information across seconds when stimuli are absent --- the cellular signature of working memory, depending on recurrent NMDA microcircuits with prominent GluN2B subunits. Joaquín Fuster generalized this to a temporal-integration function bridging perception and action across time. The decisive integration is Earl Miller and Jonathan Cohen's 2001 \emph{Annual Review} article, which proposed that \textbf{the PFC implements cognitive control by maintaining patterns of activity representing goals and the means to achieve them, which then provide biasing signals to posterior cortices and subcortical structures}. PFC is not the executor; it is the biaser of competition elsewhere in the brain.

The crucial point for the chapter's argument is that the PFC does not operate alone. It is embedded within Alexander, DeLong, and Strick's parallel cortico-striato-pallido-thalamo-cortical loops. \textbf{Each loop runs the same canonical disinhibitory selection algorithm --- the same one running in the lamprey --- applied to progressively more abstract content}: motor primitives in the sensorimotor loop, eye movements in the oculomotor loop, working-memory items and rule sets in the dorsolateral prefrontal loop, value representations in the orbitofrontal loop, and goal/effort tradeoffs in the anterior cingulate loop. Michael Frank and Randy O'Reilly's prefrontal-basal-ganglia working memory framework makes the architecture explicit: basal-ganglia Go and No-Go pathways, trained by phasic dopamine reward-prediction errors, learn when to gate information into, maintain in, or read out from PFC stripes. This decomposes the homunculus problem of executive control into a biologically plausible, reinforcement-learned gating policy. Balleine and Dickinson's dual-system view further partitions instrumental action into a \textbf{goal-directed system} (prelimbic cortex in rodents; DLPFC and caudal OFC in primates; dorsomedial striatum) sensitive to outcome devaluation, and a \textbf{habit system} (sensorimotor cortex; dorsolateral striatum) insensitive to it. With training, control shifts from goal-directed to habitual; with stress, addiction, or dopaminergic disruption, the balance breaks pathologically.

\subsection{Simulation, prediction, and the long horizon of human action}\label{05_vertebrates-simulation-prediction-and-the-long-horizon-of-human-action}

A second functional consequence of expanded PFC, in concert with the medial temporal lobe and the default mode network, is \textbf{internal simulation}. Suddendorf and Corballis introduced the concept of mental time travel as a unified faculty encompassing episodic memory and episodic foresight. Schacter, Addis, and Buckner provided the neural complement, showing that imagining future events recruits the same core network --- medial PFC, posterior cingulate and retrosplenial cortex, hippocampus, lateral parietal and lateral temporal areas --- that mediates remembering the past. Their \textbf{constructive episodic simulation hypothesis} proposes that episodic memory's flexibility is adaptive precisely because it allows the brain to simulate, predict, and pre-experience future events. The prospective brain is not a passive recorder; it is a forward-modeling engine.

This dovetails with the broader \textbf{predictive processing and active inference} framework articulated by Karl Friston and Andy Clark. The free-energy principle proposes that the brain is an inference engine that minimizes a long-term average of prediction error by either updating internal generative models (perception, learning) or acting on the world to bring sensory data into line with predictions (action). Hierarchical cortical architecture, with deeper anterior layers including PFC encoding longer temporal and more abstract regularities, instantiates the scheme. Goals become priors. Policies become inferences over likely future trajectories. Dopamine encodes the precision of prediction errors that guide policy selection. The framework unifies threads otherwise treated separately: working memory is the maintenance of high-level priors over a delay; action selection is policy inference under those priors; mental time travel is the simulation step within active inference, evaluating expected free energy of candidate policies before commitment.

What this means is that \textbf{human long-term goal pursuit --- saving for retirement, raising a child, completing a manuscript over years --- is not a different kind of operation from a lamprey selecting a swimming bout}. It is the same operation running on priors that span decades rather than seconds, scaffolded by language, culture, and external memory, and executed through the same conserved basal-ganglia gating architecture. The PFC's contribution is to make the items in competition more abstract (goals, rules, plans), to hold them across longer delays, and to nest them hierarchically so that subgoals can be selected in service of superordinate goals. The selection algorithm is unchanged.

\subsection{Conclusion: agency has a material substrate, and the substrate is older than jaws}\label{05_vertebrates-conclusion-agency-has-a-material-substrate-and-the-substrate-is-older-than-jaws}

The narrative arc of this chapter inverts the popular picture. There is no rational cortex governing an instinctual reptilian core. There is no triune brain. Cognition is not the exclusive achievement of mammalian neocortical lamination. \textbf{What there is, instead, is a 500-million-year-old vertebrate action-selection engine --- basal ganglia and dopamine --- that has been continuously elaborated across the lamprey, gnathostome, tetrapod, amniote, synapsid, mammal, and primate lineages}, with each elaboration adding capacity for richer prediction, longer time horizons, and more abstract content, while the core algorithm of disinhibition gating remained essentially unchanged. Evidence for this claim comes from comparative neuroanatomy (Striedter; Butler and Hodos), shared developmental codes (Puelles and Rubenstein), molecular cell-type homologies (Karten; Dugas-Ford and Ragsdale; Tosches), the convergent cognition of birds (Jarvis; Stacho; Güntürkün; Nieder), and most decisively from the Karolinska program in lamprey neurophysiology that has reduced the homology question to a near-identity claim at the level of cells, synapses, and circuits.

The implication for the broader manifesto this chapter serves is direct. \textbf{Human agency has a material, evolved substrate.} Whatever is involved in willing, choosing, planning, and intending is implemented --- mechanically and inspectably --- in a hierarchical control system whose deepest layer is older than the appearance of jaws. PFC-supported simulation generates candidate futures; DLPFC working-memory neurons hold a goal across a delay; basal ganglia gate that goal into action; dopamine adjusts the policy on the fly; and the whole system minimizes long-run prediction error. There is no extra ingredient --- no immaterial will, no supervisory homunculus standing outside the circuit. The smooth functioning of this stack of mechanisms is what we ordinarily call agency.

This is why the diseases of the system are diagnostic. \textbf{Parkinson's disease impairs goal-directed initiation by depleting nigrostriatal dopamine, exposing how the gating algorithm depends on dopaminergic prediction-error signals.} Addiction reflects pathological capture of the dorsolateral-striatal habit system by drug-induced dopamine surges, eroding PFC-mediated goal arbitration. Schizophrenia disrupts DLPFC delay-period firing and dopaminergic precision-weighting, fragmenting the maintenance of priors. Frontal lesions impair value-based choice and prospection. Each pathology dismantles a specific part of the very machinery whose evolutionary assembly this chapter has traced, and in doing so confirms that the machinery is what agency is \emph{made of}. The argument of subsequent chapters --- that dopamine system damage causes a breakdown of action machinery rather than a lapse of moral character --- rests on this evolutionary claim. The action machinery is real, ancient, conserved, and continuous with the rest of the vertebrate world. Agency is biology operating at its most integrated, and biology, in this domain, is older than the Cambrian closed.
